Mutation testing reliably evaluates test suite quality, but its execution cost is very high, making the selection of a useful mutant subset an important practical problem. Prior studies have used LLMs to generate mutants or detect equivalent mutants, but none have evaluated an LLM as a mutant-selection decision maker against random selection under the same budget. We use GPT-4o-mini (temperature=0, fixed prompt) to rank 11,040 PIT mutants from three Defects4J projects (Commons Codec, Commons CSV, JacksonCore; versions 1f–10f), select the top 10\% within each (project, version, class) group, and compare it with random selection under the same budget (seed 42) using mutation score. GPT-4o-mini achieves an overall mutation score of 0.716 compared to 0.670 for random selection (+4.6 points); a one-sided Wilcoxon test at the group level gives p=0.059, r\_rb=+0.350 (not meeting the full pre-registered criteria), while mutant-weighted analysis (p=0.0097) and a 5,000-draw Monte Carlo simulation (p=0.0002) support GPT's advantage. These results indicate that LLM-based mutant selection provides concentrated benefits on large, complex parser classes but does not consistently outperform random selection across all groups — consistent with earlier findings on naturalness-based selection.
